Pellen iemand van vooren
Veronderstelling / voorspel

voorst
meest vooraan, eerder dan alle andere
Het neologisme "voorstpelling" lijkt een samensmelting van de woorden "voorstelling" (het verbeelden, tonen van iets) en "voorspelling" (het vooraf zeggen wat zal gebeuren). De betekenis van "voorstpelling" kan daarom als volgt worden uitgelegd:
Uitleg van het neologisme "voorstpelling"
Voorstpelling verwijst naar een imaginair of creatief proces waarin het toekomstige (dat wat te voorspellen valt) tegelijk wordt voorgesteld en voorspeld. Het is dus niet enkel het rationeel voorzien van de toekomst (voorspelling), noch puur het artistiek verbeelden ervan (voorstelling), maar een tussenstadium:
Een voorstpelling maakt het onzekere, het toekomstige zichtbaar binnen een (kunstzinnige) verbeelding, en verleent er daardoor een voorlopige werkelijkheid aan. Het kan gebruikt worden om het proces te beschrijven waarin literatuur, muziek of kunst een beeld geeft van hoe iets zal zijn, zonder dat het strikt feitelijk is—een voorstelbare voorspelling.
Mogelijke toepassingen binnen Emmy Nativus
In de context van Emmy Nativus kan "voorstpelling" bijvoorbeeld verwijzen naar het proces waarbij de schrijver niet slechts zijn innerlijke leven voorspelt (analyseert, rationeel benadert), maar ook voorstelt (verbeeldt, artistiek vormgeeft)—en deze twee processen in één beweging samenbrengt.
Samenvatting
Voorstpelling = Voorspelling + Voorstelling
Betekenis: Een verbeelde voorspelling; het voorstelbaar maken van het toekomstige door kunst, muziek of literatuur, waarbij feit en creatie samenvloeien.
Baby laat sleutels vallen zodat brein voorspel genot wordt bevredigd totdat sleutels ineens omhoog gaan. Dan angst door onzekerheid ná eerst de verwondering. Daarom laat de baby het vallen omdat het brein dan kan laten zien dat het werkt steeds weer en dat geeft een veilig gevoel, ik kan de wereld voorspellen en dus controleren.
